\documentclass[10pt, aspectratio=169]{beamer}
\usetheme{Madrid}
\usecolortheme{dolphin}
\usepackage{ctex}
\usepackage{xeCJK}
\usepackage{amsmath,amssymb}
\usepackage{listings}
\usepackage{tikz}
\usepackage{xcolor}
\title{408考研零基础 --- 数据结构}
\subtitle{1-17 查找基础}
\author{}
\date{}
\setbeamertemplate{page number in head/foot}{}
\begin{document}
\frame{\titlepage}
\begin{frame}{本节目标}
\begin{enumerate}
\item 掌握顺序查找的算法与性能
\item 掌握折半查找的算法与判定树
\item 理解分块查找和 B 树基本概念
\end{enumerate}
\end{frame}
\begin{frame}{查找的基本概念}
\textbf{平均查找长度 ASL}
\[
\text{ASL} = \sum_{i=1}^{n} p_i \times c_i
\]
\vspace{0.3em}
\textbf{分类}
\begin{itemize}
\item 静态查找表：只查找，不修改
\item 动态查找表：查找同时插入/删除
\end{itemize}
\end{frame}
\begin{frame}{顺序查找}
\textbf{算法}：从表一端开始依次比较
\vspace{0.3em}
\textbf{性能}
\begin{itemize}
\item 查找成功 ASL：$(n+1)/2$
\item 查找失败比较次数：$n+1$
\item 时间复杂度 $O(n)$
\end{itemize}
\vspace{0.3em}
\textbf{优化}：设置哨兵减少边界判断
\vspace{0.3em}
适用于无序表或小规模数据
\end{frame}
\begin{frame}{折半查找}
\textbf{适用条件}：有序的顺序表
\vspace{0.3em}
\textbf{步骤}
\begin{itemize}
\item low, high 指向两端，mid $=$ (low+high)/2
\item 关键字 $=$ mid 元素 $\to$ 成功
\item 关键字 $<$ mid $\to$ high $=$ mid$-1$
\item 关键字 $>$ mid $\to$ low $=$ mid$+1$
\item low $>$ high $\to$ 失败
\end{itemize}
\vspace{0.3em}
\textbf{性能}：ASL $\approx \log_2(n+1)-1$，$O(\log n)$
\end{frame}
\begin{frame}{分块查找}
\textbf{块间有序，块内无序}
\vspace{0.3em}
\textbf{步骤}
\begin{enumerate}
\item 索引表中折半查找确定所属块
\item 块内顺序查找
\end{enumerate}
\vspace{0.3em}
\textbf{性能}：$s=\sqrt{n}$ 时 ASL 最小 $\approx \sqrt{n}+1$
\vspace{0.3em}
适合动态插入场景
\end{frame}
\begin{frame}{B 树的基本概念}
$m$ 阶 B 树
\begin{itemize}
\item 每个结点至多 $m$ 棵子树（$m-1$ 个关键字）
\item 根至少 2 棵子树
\item 非叶结点至少 $\lceil m/2 \rceil$ 棵子树
\item 所有叶结点在同一层
\end{itemize}
\vspace{0.3em}
\textbf{查找}：$O(\log_m n)$，降低树高减少磁盘 I/O
\vspace{0.3em}
广泛应用于文件系统和数据库索引
\end{frame}
\begin{frame}{本节总结}
\begin{enumerate}
\item 顺序查找：$O(n)$，ASL $(n+1)/2$
\item 折半查找：$O(\log n)$，要求有序顺序表
\item 分块查找：索引 + 顺序，适合动态场景
\item B 树：多路平衡查找，$O(\log_m n)$，用于数据库索引
\end{enumerate}
\vspace{1em}
\begin{center}
\textcolor{blue}{\textbf{下节预告：排序入门}}
\end{center}
\end{frame}
\end{document}
